% =====================================================================
%  CHƯƠNG 2 — CÁC KỸ THUẬT HỌC MÁY CƠ BẢN
% =====================================================================
\chapter{Các Kỹ Thuật Học Máy Cơ Bản}
\label{ch:hoc-may}

\section{Cơ sở lý thuyết các thuật toán học máy cốt lõi}

Học máy có giám sát giải hai lớp bài toán: \emph{hồi quy} (dự báo một giá trị liên tục) và \emph{phân loại} (gán
nhãn rời rạc). Cho tập huấn luyện $\{(\mathbf{x}^{(i)}, y^{(i)})\}_{i=1}^{N}$ với $\mathbf{x} \in \mathbb{R}^{D}$, ta chọn
một họ hàm $f_\theta$ và tìm tham số $\theta$ cực tiểu hóa một hàm mất mát trung bình trên dữ liệu. Bốn họ thuật toán
dưới đây khác nhau ở họ hàm và cách tối ưu \cite{raschka2015,goodfellow2016}.

\subsection{Hồi quy tuyến tính (Linear Regression)}

Mô hình giả định đầu ra là tổ hợp tuyến tính của đặc trưng:
\begin{equation}
  \hat y = f(\mathbf{x}) = \mathbf{w}^\top\mathbf{x} + b = \sum_{j=1}^{D} w_j x_j + b ,
\end{equation}
và cực tiểu hóa sai số toàn phương trung bình
\begin{equation}
  \mathcal{L}_{\mathrm{MSE}}(\mathbf{w}, b) = \frac{1}{2N}\sum_{i=1}^{N}\big(\hat y^{(i)} - y^{(i)}\big)^2
    = \frac{1}{2N}\lVert X\mathbf{w} + b\mathbf{1} - \mathbf{y}\rVert_2^2 .
\end{equation}
Nghiệm đóng là phương trình chuẩn $\mathbf{w}^* = (X^\top X)^{-1}X^\top\mathbf{y}$; khi $D$ lớn có thể dùng hạ gradient
$\mathbf{w} \leftarrow \mathbf{w} - \eta\,\tfrac{1}{N}X^\top(X\mathbf{w} - \mathbf{y})$. Khi các đặc trưng tương quan mạnh
(đa cộng tuyến), $X^\top X$ gần suy biến và hệ số dao động lớn; \textbf{Ridge} thêm phạt $\alpha\lVert\mathbf{w}\rVert_2^2$
để nghiệm thành $(X^\top X + \alpha I)^{-1}X^\top\mathbf{y}$, luôn khả nghịch.

\subsection{Hồi quy Logistic (Logistic Regression)}

Với nhãn $y \in \{0, 1\}$, tổ hợp tuyến tính được ép về khoảng $(0, 1)$ bằng hàm sigmoid:
\begin{equation}
  p = P(y = 1 \mid \mathbf{x}) = \sigma(z) = \frac{1}{1 + e^{-z}}, \qquad z = \mathbf{w}^\top\mathbf{x} + b .
\end{equation}
Ước lượng hợp lý cực đại dẫn tới hàm mất mát entropy chéo nhị phân
\begin{equation}
  \mathcal{L}_{\mathrm{BCE}} = -\frac{1}{N}\sum_{i=1}^{N}\Big[y^{(i)}\log\hat y^{(i)} + (1 - y^{(i)})\log(1 - \hat y^{(i)})\Big],
  \label{eq:bce}
\end{equation}
có đạo hàm theo $z$ rút gọn thành $\partial\mathcal{L}/\partial z = \hat y - y$, nên
$\nabla_{\mathbf{w}}\mathcal{L} = \tfrac{1}{N}\sum_i (\hat y^{(i)} - y^{(i)})\mathbf{x}^{(i)}$. Mất mát là hàm lồi, mọi
cực tiểu đều là toàn cục. Với $K$ lớp, sigmoid được thay bằng softmax $\hat y_k = e^{z_k}/\sum_j e^{z_j}$ và mất mát là
entropy chéo $-\log \hat y_{y}$. Khi dữ liệu mất cân bằng, trọng số lớp \texttt{class\_weight="balanced"} nhân mất mát của
mỗi mẫu với $N/(K\,N_k)$ để lớp hiếm có tiếng nói ngang lớp phổ biến.

\subsection{Cây quyết định và Rừng ngẫu nhiên (Random Forest)}

Cây quyết định chia không gian bằng chuỗi câu hỏi nhị phân “$x_j \le \theta$?”. Độ thuần của một nút $t$ có các lớp với tỉ
lệ $p_k$ được đo bằng chỉ số Gini hoặc entropy:
\begin{equation}
  \mathrm{Gini}(t) = 1 - \sum_{k=1}^{K} p_k^2, \qquad \mathrm{Entropy}(t) = -\sum_{k=1}^{K} p_k\log_2 p_k .
\end{equation}
Tại mỗi nút, thuật toán chọn cặp $(j, \theta)$ làm giảm tạp chất nhiều nhất:
\begin{equation}
  \Delta\mathrm{Gini} = \mathrm{Gini}(t) - \Big(\frac{N_L}{N_t}\mathrm{Gini}(t_L) + \frac{N_R}{N_t}\mathrm{Gini}(t_R)\Big).
\end{equation}
Với hồi quy, tạp chất là phương sai của nhãn trong nút. Một cây sâu khớp rất sát dữ liệu huấn luyện nhưng có phương sai
lớn. \textbf{Random Forest} \cite{breiman2001} trung bình hóa $B$ cây với hai nguồn ngẫu nhiên: mỗi cây học trên một mẫu
bootstrap (rút có hoàn lại $N$ mẫu) và mỗi lần chia chỉ xét một tập con ngẫu nhiên $\sqrt{D}$ đặc trưng. Hai cơ chế này
làm các cây ít tương quan với nhau, nên trung bình của chúng có phương sai nhỏ hơn hẳn từng cây mà độ chệch gần như giữ
nguyên.

\subsection{Máy vector hỗ trợ (Support Vector Machine)}

SVM \cite{cortes1995} tìm siêu phẳng phân tách có lề lớn nhất, cho phép vi phạm qua biến nới lỏng $\xi_i$:
\begin{equation}
  \min_{\mathbf{w}, b, \boldsymbol{\xi}}\ \frac{1}{2}\lVert\mathbf{w}\rVert^2 + C\sum_{i=1}^{N}\xi_i
  \quad\text{với}\quad y^{(i)}(\mathbf{w}^\top\mathbf{x}^{(i)} + b) \ge 1 - \xi_i,\ \ \xi_i \ge 0 ,
\end{equation}
trong đó $C$ cân bằng giữa lề rộng và lỗi huấn luyện. Bài toán đối ngẫu chỉ phụ thuộc vào tích vô hướng giữa các mẫu, nên
có thể thay bằng một hàm nhân. Nhân RBF
\begin{equation}
  K(\mathbf{x}, \mathbf{x}') = \exp\!\big(-\gamma\lVert\mathbf{x} - \mathbf{x}'\rVert^2\big)
\end{equation}
tương ứng với ánh xạ lên không gian vô hạn chiều. Cái giá là bộ giải phải làm việc với ma trận nhân $N \times N$: thời gian
huấn luyện tăng cỡ $O(N^2)$ tới $O(N^3)$ và mô hình phải lưu toàn bộ vector hỗ trợ. \textbf{SVR} là phiên bản hồi quy, phạt
các điểm nằm ngoài ống sai số rộng $\varepsilon$. Với văn bản TF-IDF hàng nghìn chiều, dữ liệu thường đã gần khả tách tuyến
tính nên \textbf{LinearSVC} (nhân tuyến tính, độ phức tạp tuyến tính theo $N$) là lựa chọn hợp lý hơn.

% ---------------------------------------------------------------------
\section{Khảo sát thực nghiệm và tiền xử lý ba bộ dữ liệu}

Ba tập dữ liệu được chọn để phủ ba dạng bài toán: phân loại nhị phân mất cân bằng (y tế), hồi quy (bất động sản Việt Nam)
và phân loại văn bản ba lớp (thương mại điện tử). Mọi tập đều chia 80\,\% huấn luyện / 20\,\% kiểm thử với
\texttt{random\_state = 42} (phân tầng theo nhãn với bài toán phân loại); mạng MLP tách thêm 10\,\% của phần huấn luyện làm
tập kiểm định cho dừng sớm.

\subsection{Tập dữ liệu 1: Sàng lọc nguy cơ tiểu đường (\SL{eda/diabetes/n} mẫu)}

\begin{itemize}
  \item \textbf{Nguồn:} Diabetes Prediction Dataset trên Kaggle \cite{kaggle_diabetes}, tổng hợp hồ sơ khám bệnh.
  \item \textbf{Quy mô:} \SL{eda/diabetes/n} hồ sơ, \SL{eda/diabetes/n_features} thuộc tính: giới tính, tuổi, tăng huyết áp,
    bệnh tim, tiền sử hút thuốc, chỉ số khối cơ thể BMI, HbA1c (\%), đường huyết (mg/dL); nhãn \texttt{diabetes}
    (0: không mắc, 1: mắc).
\end{itemize}

\bang{tab_sample_diab}{Năm bản ghi ngẫu nhiên của tập dữ liệu tiểu đường.}{sample-diab}

\hinhpdf{f2_diab_target}{0.62}{Phân bố nhãn của tập tiểu đường: lớp dương chỉ chiếm \SL{eda/diabetes/pos_ratePct}\,\%.}{diab-target}
\hinhpdf{f2_diab_kde}{1}{Mật độ xác suất của đường huyết, HbA1c, BMI và tuổi theo tình trạng bệnh.}{diab-kde}
\hinhpdf{f2_diab_corr}{0.72}{Ma trận tương quan Pearson giữa các chỉ số lâm sàng và nhãn tiểu đường.}{diab-corr}

\paragraph{Nhận xét EDA.}
\begin{itemize}
  \item \textbf{Mất cân bằng lớp:} chỉ \SL{eda/diabetes/n_pos} trên \SL{eda/diabetes/n} hồ sơ là dương tính
    (Hình~\ref{fig:diab-target}). Một mô hình ngây thơ luôn đoán “không mắc” đã đạt Accuracy
    $100 - \SL{eda/diabetes/pos_ratePct} \approx 91{,}5\,\%$ mà không phát hiện được ca bệnh nào; vì vậy Accuracy không
    đủ để đánh giá, cần Precision, Recall và F1 của lớp dương cùng ROC-AUC.
  \item \textbf{Hai đặc trưng quyết định:} đường huyết và HbA1c có tương quan với nhãn lần lượt
    $r = \SL{eda/diabetes/corr_glucose}$ và $r = \SL{eda/diabetes/corr_hba1c}$, cao nhất trong các chỉ số
    (Hình~\ref{fig:diab-corr}). Phân bố của người mắc bệnh dịch hẳn sang phải (Hình~\ref{fig:diab-kde}): tỉ lệ mắc bệnh là
    \SL{eda/diabetes/rate_glu200Pct}\,\% khi đường huyết $\ge 200$\,mg/dL, \SL{eda/diabetes/rate_hba65Pct}\,\% khi
    HbA1c $\ge 6{,}6$\,\%, nhưng chỉ \SL{eda/diabetes/rate_hba_lowPct}\,\% khi HbA1c $< 5{,}7$\,\%.
  \item \textbf{Quan hệ dạng ngưỡng:} cả hai chỉ số đều có “ngưỡng gãy” rõ (đặc biệt HbA1c), một dạng phi tuyến mà
    mô hình tuyến tính không biểu diễn trọn được nhưng cây quyết định và mạng nơ-ron thì có thể.
  \item \textbf{Tiền xử lý:} chuẩn hóa z-score bốn biến số, giữ nguyên hai biến nhị phân, mã hóa one-hot giới tính và tiền
    sử hút thuốc (\SL{eda/diabetes/noinfo_smokePct}\,\% hồ sơ ghi “No Info”, được giữ thành một nhóm riêng).
\end{itemize}

\subsection{Tập dữ liệu 2: Định giá bất động sản Việt Nam (\SL{eda/house/n} mẫu)}

\begin{itemize}
  \item \textbf{Nguồn:} Vietnam Housing Dataset \cite{kaggle_vnhouse}, tin đăng bán nhà trên các trang bất động sản Việt Nam.
  \item \textbf{Quy mô:} \SL{eda/house/n} tin đăng thuộc \SL{eda/house/n_province} tỉnh/thành và \SL{eda/house/n_district}
    quận/huyện; các thuộc tính diện tích, mặt tiền, độ rộng đường vào, số tầng, số phòng ngủ, số phòng tắm, hướng nhà,
    hướng ban công, tình trạng pháp lý, nội thất, địa chỉ; nhãn là giá bán (tỷ đồng).
\end{itemize}

\bang{tab_sample_house}{Năm tin đăng ngẫu nhiên của tập dữ liệu giá nhà.}{sample-house}

\hinhpdf{f2_house_dist}{1}{Phân bố giá bán, diện tích gốc và diện tích sau biến đổi $\log(1 + x)$.}{house-dist}
\hinhpdf{f2_house_scatter}{1}{Trái: giá theo diện tích (thang log), tô màu theo độ rộng đường vào. Phải: giá theo sáu
tỉnh/thành có nhiều tin nhất.}{house-scatter}
\hinhpdf{f2_house_corr}{0.62}{Ma trận tương quan giữa các thuộc tính vật lý và giá bán.}{house-corr}

\paragraph{Nhận xét EDA.}
\begin{itemize}
  \item \textbf{Giá đã bị chặn:} giá nằm trong $[\SL{eda/house/price_min}; \SL{eda/house/price_max}]$ tỷ đồng, trung bình
    \SL{eda/house/price_mean}, độ lệch $\SL{eda/house/price_skew}$ (Hình~\ref{fig:house-dist}). Phân bố gần đối xứng nên
    \emph{không} cần biến đổi log cho nhãn; ngược lại diện tích lệch phải mạnh (độ lệch \SL{eda/house/area_skew}), sau
    $\log(1 + x)$ còn \SL{eda/house/logarea_skew}.
  \item \textbf{Tương quan yếu và phi tuyến:} tương quan tuyến tính của giá với diện tích chỉ $r = \SL{eda/house/corr_area}$,
    với số tầng $r = \SL{eda/house/corr_floors}$, với đường vào $r = \SL{eda/house/corr_road}$ (Hình~\ref{fig:house-corr}).
    Giá phụ thuộc mạnh vào vị trí: trung vị giá ở TP.~Hồ Chí Minh và Hà Nội cao nhất, Đà Nẵng thấp hơn một chút, còn Bình
    Dương và Đồng Nai thấp hơn gần một nửa (Hình~\ref{fig:house-scatter}), mà vị trí là biến phân loại có hàng trăm giá trị.
  \item \textbf{Thiếu dữ liệu nhiều:} \SL{eda/house/miss_frontagePct}\,\% tin thiếu mặt tiền, \SL{eda/house/miss_roadPct}\,\%
    thiếu đường vào, \SL{eda/house/miss_directionPct}\,\% thiếu hướng nhà.
  \item \textbf{Tiền xử lý:} tách tỉnh và quận/huyện từ địa chỉ; biến số được điền trung vị \emph{kèm cờ đánh dấu thiếu} rồi
    chuẩn hóa z-score; biến phân loại điền “Không rõ” rồi mã hóa one-hot, gộp các giá trị xuất hiện dưới 40 lần thành một
    nhóm “hiếm” để không sinh hàng trăm cột thưa.
\end{itemize}

\subsection{Tập dữ liệu 3: Phân loại cảm xúc nhận xét thương mại điện tử (\SL{eda/reviews/n} mẫu)}

\begin{itemize}
  \item \textbf{Nguồn:} Women's E-Commerce Clothing Reviews \cite{kaggle_reviews}, nhận xét thật (tiếng Anh) của khách
    hàng một nhà bán lẻ thời trang trực tuyến.
  \item \textbf{Quy mô:} \SL{eda/reviews/n_raw} bản ghi, giữ \SL{eda/reviews/n} bản có nội dung nhận xét. Nhãn cảm xúc ba
    lớp suy từ số sao: 1 -- 2 sao là \emph{tiêu cực}, 3 sao là \emph{trung tính}, 4 -- 5 sao là \emph{tích cực}. Tiêu đề và
    nội dung được ghép thành một văn bản.
\end{itemize}

\bang{tab_sample_rev}{Năm nhận xét tiêu biểu kèm nhãn cảm xúc.}{sample-rev}

\hinhpdf{f2_rev_dist}{0.95}{Phân bố số sao và nhãn cảm xúc: lớp tích cực chiếm \SL{eda/reviews/share_posPct}\,\%.}{rev-dist}
\hinhpdf{f2_rev_len}{0.65}{Phân bố độ dài nhận xét (số từ) theo nhãn; nét đứt là trung bình \SL{eda/reviews/words_mean} từ.}{rev-len}
\hinhpdf{f2_rev_words}{1}{Mười lăm từ đặc trưng nhất của nhận xét tích cực và tiêu cực, xếp theo tỉ số log-odds so với
phần còn lại.}{rev-words}

\paragraph{Nhận xét EDA và quy trình tiền xử lý văn bản.} Dữ liệu lệch mạnh về phía tích cực
(\SL{eda/reviews/share_posPct}\,\%), trong khi trung tính chỉ \SL{eda/reviews/share_neuPct}\,\% và tiêu cực
\SL{eda/reviews/share_negPct}\,\% (Hình~\ref{fig:rev-dist}); vì vậy chỉ số chính là F1 trung bình macro trên ba lớp.
Độ dài gần như không phân biệt được các lớp: trung bình \SL{eda/reviews/words_by_class/0} từ ở nhận xét tiêu cực,
\SL{eda/reviews/words_by_class/1} ở trung tính và \SL{eda/reviews/words_by_class/2} ở tích cực (Hình~\ref{fig:rev-len}),
nên tín hiệu nằm ở \emph{từ ngữ} chứ không ở độ dài. Từ đặc trưng nhất của lớp tiêu cực là \emph{\SL{eda/reviews/top_neg/0}},
\emph{\SL{eda/reviews/top_neg/1}}, \emph{\SL{eda/reviews/top_neg/2}}, \emph{\SL{eda/reviews/top_neg/3}}; của lớp tích cực là
\emph{\SL{eda/reviews/top_pos/3}}, \emph{\SL{eda/reviews/top_pos/4}}, \emph{\SL{eda/reviews/top_pos/0}}
(Hình~\ref{fig:rev-words}). Quy trình tiền xử lý:
\begin{enumerate}
  \item chuyển về chữ thường, bỏ URL và mọi ký tự không phải chữ cái hay dấu nháy;
  \item loại từ dừng tiếng Anh (danh sách của scikit-learn);
  \item vector hóa TF-IDF với unigram và bigram, giữ 5.000 đặc trưng tần suất cao nhất, bỏ từ xuất hiện dưới 3 văn bản,
    dùng $1 + \log(\text{tf})$ để giảm ảnh hưởng của từ lặp nhiều lần.
\end{enumerate}

% ---------------------------------------------------------------------
\section{Hiện thực hóa và thực nghiệm các mô hình học máy cơ bản}

Mỗi mô hình được đóng gói thành một \texttt{Pipeline} của scikit-learn \cite{pedregosa2011} gồm bộ tiền xử lý và bộ ước
lượng. Gọi \texttt{fit} trên phần huấn luyện sẽ khớp cả hai, nên không thống kê nào của tập kiểm thử rò vào quá trình
học; lưu pipeline bằng \texttt{joblib} là có ngay đối tượng nhận dữ liệu thô cho Web App.

\subsection{Pipeline hồi quy tuyến tính và Ridge}

Bộ tiền xử lý giá nhà (Đoạn mã~\ref{code:data_tabular__house_preprocessor}) xử lý riêng nhóm cột số và nhóm cột phân loại
bằng \texttt{ColumnTransformer}; Đoạn mã~\ref{code:ml_classic__build_linear_regression_pipeline} ghép nó với mô hình
tuyến tính, \texttt{alpha=None} cho bình phương tối thiểu thường và \texttt{alpha=1.0} cho Ridge.

\maNguon{data_tabular__house_preprocessor}{Bộ tiền xử lý dữ liệu giá nhà: điền thiếu kèm cờ, chuẩn hóa, one-hot có gộp nhóm hiếm}
\maNguon{ml_classic__build_linear_regression_pipeline}{Pipeline hồi quy tuyến tính và Ridge với scikit-learn}

\subsection{Pipeline hồi quy Logistic cho dữ liệu y tế}

\texttt{class\_weight="balanced"} bù cho tỉ lệ dương \SL{eda/diabetes/pos_ratePct}\,\%; bộ giải \texttt{lbfgs} là phương
pháp quasi-Newton hội tụ nhanh với số đặc trưng nhỏ.

\maNguon{ml_classic__build_logistic_regression_pipeline}{Pipeline hồi quy Logistic có cân bằng trọng số lớp}

\subsection{Pipeline Rừng ngẫu nhiên}

Rừng gồm 200 cây không giới hạn độ sâu, mỗi lá tối thiểu 2 mẫu; với phân loại, trọng số lớp được cân bằng lại trên từng
mẫu bootstrap. Cây không cần chuẩn hóa đặc trưng, nhưng dùng chung bộ tiền xử lý giúp mọi mô hình nhận cùng một đầu vào.

\maNguon{ml_classic__build_rf_pipeline}{Pipeline Random Forest cho cả hồi quy và phân loại}

\subsection{Pipeline Máy vector hỗ trợ}

Nhân RBF tính khoảng cách Euclid nên đặc trưng bắt buộc cùng thang đo; bộ tiền xử lý đã chuẩn hóa z-score. Văn bản dùng
\texttt{LinearSVC}.

\maNguon{ml_classic__build_svm_pipeline}{Pipeline SVM: SVC/SVR nhân RBF cho dữ liệu bảng, LinearSVC cho văn bản}

\subsection{Kết quả thực nghiệm các mô hình học máy cơ bản}

Ba bảng dưới đây tách theo bản chất bài toán. Thời gian là thời gian huấn luyện trên CPU (Random Forest dùng 12 luồng);
dung lượng là kích thước tệp \texttt{joblib} nén.

\bang{tab_ch2_diab_sk}{Kết quả các mô hình học máy cơ bản trên tập tiểu đường (\SL{eda/diabetes/n} hồ sơ, kiểm thử 20\,\%).}{ch2-diab-sk}
\bang{tab_ch2_house_sk}{Kết quả các mô hình học máy cơ bản trên tập giá nhà Việt Nam.}{ch2-house-sk}
\bang{tab_ch2_rev_sk}{Kết quả các mô hình học máy cơ bản trên tập nhận xét khách hàng (TF-IDF 5.000 chiều).}{ch2-rev-sk}

\subsection{Phân tích kết quả và giới hạn của các mô hình cơ bản}

\begin{enumerate}
  \item \textbf{Mô hình tuyến tính.} Trên tập tiểu đường, Logistic Regression có trọng số cân bằng đạt Recall
    \SL{ch2_sklearn/diabetes/logreg/recall} nhưng Precision chỉ \SL{ch2_sklearn/diabetes/logreg/precision}: trọng số lớp
    đẩy ngưỡng quyết định về phía bắt nhiều ca bệnh, đổi lại hơn một nửa cảnh báo là báo động giả. ROC-AUC
    \SL{ch2_sklearn/diabetes/logreg/auc} cho thấy khả năng \emph{xếp hạng} vẫn tốt; điểm yếu nằm ở ranh giới tuyến tính
    không bắt được ngưỡng gãy của HbA1c. Trên tập giá nhà, OLS và Ridge cùng chỉ đạt $R^2 \approx \SL{ch2_sklearn/house/ols/r2}$
    (MAE \SL{ch2_sklearn/house/ols/mae} tỷ): phạt L2 gần như không giúp gì vì vấn đề là thiếu khớp chứ không phải đa
    cộng tuyến.
  \item \textbf{SVM nhân RBF.} SVC nâng ROC-AUC lên \SL{ch2_sklearn/diabetes/svm/auc} và F1 lên
    \SL{ch2_sklearn/diabetes/svm/f1}, nhưng huấn luyện mất \SL{ch2_sklearn/diabetes/svm/train_s}\,s trên
    80.000 mẫu, gấp hàng trăm lần Logistic Regression, và suy luận chậm nhất (\SL{ch2_sklearn/diabetes/svm/infer_ms}\,ms
    mỗi mẫu) vì phải so mẫu mới với mọi vector hỗ trợ. SVR đạt $R^2 = \SL{ch2_sklearn/house/svr/r2}$ trên giá nhà. Với văn
    bản, LinearSVC huấn luyện trong \SL{ch2_sklearn/reviews/linsvm/train_s}\,s và đạt Accuracy
    \SL{ch2_sklearn/reviews/linsvm/accPct}\,\%.
  \item \textbf{Random Forest.} Là mô hình cơ bản mạnh nhất trên cả hai tập bảng: F1 \SL{ch2_sklearn/diabetes/rf/f1} trên
    tiểu đường và $R^2 = \SL{ch2_sklearn/house/rf/r2}$, MAE \SL{ch2_sklearn/house/rf/mae} tỷ trên giá nhà. Đổi lại, tệp mô
    hình nặng \SL{ch2_sklearn/house/rf/size_kb}\,KB (200 cây sâu), không ngoại suy được ra ngoài khoảng giá đã thấy, và trên
    văn bản 5.000 chiều thưa thì kém hơn mô hình tuyến tính (F1 macro \SL{ch2_sklearn/reviews/rf/f1}) vì mỗi lần chia
    chỉ xét một từ.
  \item \textbf{Động lực chuyển sang MLP.} Mô hình tuyến tính thiếu khả năng biểu diễn, SVM nhân không mở rộng theo số
    mẫu, Random Forest cồng kềnh. Mạng nơ-ron đa tầng hứa hẹn cả ba: xấp xỉ được hàm phi tuyến bất kỳ (định lý xấp xỉ phổ
    quát), huấn luyện theo lô nhỏ nên chi phí tuyến tính theo $N$ và chạy song song trên GPU, tệp mô hình chỉ vài chục
    đến vài trăm KB.
\end{enumerate}

% ---------------------------------------------------------------------
\section{Mạng nơ-ron đa tầng (MLP) trên ba nền tảng: NumPy, Keras, PyTorch}

\subsection{Cơ sở lý thuyết và kiến trúc mạng}

Một MLP gồm các tầng kết nối đầy đủ. Tầng $l$ nhận kích hoạt $\mathbf{a}^{[l-1]}$ của tầng trước và tính
\begin{equation}
  \mathbf{z}^{[l]} = W^{[l]}\mathbf{a}^{[l-1]} + \mathbf{b}^{[l]}, \qquad \mathbf{a}^{[l]} = g^{[l]}\big(\mathbf{z}^{[l]}\big),
\end{equation}
với $g$ là ReLU ở tầng ẩn và sigmoid, softmax hoặc hàm đồng nhất ở tầng ra tùy bài toán. Lan truyền ngược áp dụng quy tắc
dây chuyền từ tầng ra về tầng vào:
\begin{align}
  \boldsymbol{\delta}^{[L]} &= \hat{\mathbf{y}} - \mathbf{y}, \\
  \boldsymbol{\delta}^{[l]} &= \big((W^{[l+1]})^\top\boldsymbol{\delta}^{[l+1]}\big) \odot g'^{[l]}\big(\mathbf{z}^{[l]}\big), \\
  \frac{\partial\mathcal{L}}{\partial W^{[l]}} &= \boldsymbol{\delta}^{[l]}\big(\mathbf{a}^{[l-1]}\big)^\top, \qquad
  \frac{\partial\mathcal{L}}{\partial\mathbf{b}^{[l]}} = \boldsymbol{\delta}^{[l]} .
\end{align}
Đẳng thức đầu đúng cho cả ba cặp đầu ra -- mất mát dùng ở đây (sigmoid + BCE, softmax + entropy chéo, tuyến tính + MSE
với hệ số 2), nên một hàm \texttt{backward} phục vụ cả ba bài toán. Trọng số khởi tạo theo He \cite{he2015init}
($\mathcal{N}(0, 2/n_{in})$) để biên độ kích hoạt ổn định qua các tầng ReLU, và cập nhật bằng Adam \cite{kingma2015}.

Để so sánh ba nền tảng một cách công bằng, cả ba cài đặt dùng \emph{cùng một kiến trúc} (Hình~\ref{fig:mlp}): hai tầng ẩn
128 và 64 nơ-ron, ReLU, Dropout 0,2 \cite{srivastava2014}; Adam với tốc độ học $10^{-3}$, lô 256; dừng sớm khi mất mát
kiểm định không giảm trong 8 epoch và khôi phục trọng số tốt nhất.

\hinhtikz{f2_mlp}{Kiến trúc MLP dùng chung cho cả ba nền tảng và ba bài toán.}{mlp}

\subsection{Hiện thực từ đầu bằng NumPy}

Lớp \texttt{MLPScratch} chỉ dùng phép toán ma trận của NumPy. Lan truyền tiến (Đoạn mã~\ref{code:mlp_scratch__MLPScratch.forward})
lưu lại $\mathbf{z}$, $\mathbf{a}$ và mặt nạ Dropout của từng tầng; lan truyền ngược
(Đoạn mã~\ref{code:mlp_scratch__MLPScratch.backward}) dùng lại chúng để tính gradient đúng như các công thức trên; Adam
(Đoạn mã~\ref{code:mlp_scratch__MLPScratch._adam}) giữ hai mômen cho mỗi tham số và hiệu chỉnh độ chệch khởi đầu.

\maNguon{mlp_scratch__MLPScratch.forward}{Lan truyền tiến của MLP viết tay bằng NumPy, kèm Dropout ngược}
\maNguon{mlp_scratch__MLPScratch.backward}{Lan truyền ngược của MLP viết tay: một sai số tầng ra chung cho ba bài toán}
\maNguon{mlp_scratch__MLPScratch._adam}{Bộ tối ưu Adam viết tay}

\subsection{Hiện thực bằng Keras / TensorFlow}

Keras gói cả kiến trúc, hàm mất mát và vòng huấn luyện vào API cấp cao; phần đầu ra chọn theo bảng \texttt{HEADS}.

\maNguon{mlp_keras__build_keras_mlp}{MLP bằng Keras Sequential API cho ba loại đầu ra}

\subsection{Hiện thực bằng PyTorch}

PyTorch tách mô hình (\texttt{nn.Module}) khỏi vòng huấn luyện do ta tự viết. Mô hình trả logit thô; sigmoid và softmax
nằm trong \texttt{BCEWithLogitsLoss} và \texttt{CrossEntropyLoss} để ổn định số học. Vì dữ liệu bảng nhỏ, toàn bộ tập
được đưa lên GPU một lần và chia lô bằng chỉ số thay vì \texttt{DataLoader}.

\maNguon{mlp_torch__TorchMLP}{MLP bằng PyTorch với lớp \texttt{nn.Module}}
\maNguon{mlp_torch__train_epoch}{Một epoch huấn luyện PyTorch trên dữ liệu đã nằm sẵn trên GPU}

% ---------------------------------------------------------------------
\section{Đánh giá và đối chuẩn trên ba tập dữ liệu}

\subsection{Đối chuẩn hiệu năng chi tiết trên từng tập dữ liệu}

Ba bảng sau đặt ba bản MLP cạnh các mô hình cơ bản, trên cùng tập kiểm thử.

\bang{tab_ch2_diab_all}{Đối chuẩn mọi mô hình trên tập tiểu đường.}{ch2-diab-all}
\bang{tab_ch2_house_all}{Đối chuẩn mọi mô hình trên tập giá nhà Việt Nam.}{ch2-house-all}
\bang{tab_ch2_rev_all}{Đối chuẩn mọi mô hình trên tập nhận xét khách hàng.}{ch2-rev-all}
\bang{tab_ch2_fw}{Cùng một kiến trúc MLP trên ba nền tảng.}{ch2-fw}

\hinhpdf{f2_loss_mlp}{1}{Đường cong mất mát huấn luyện (nét đứt) và kiểm định (nét liền) của MLP trên ba nền tảng.}{loss-mlp}
\hinhpdf{f2_house_pred}{0.95}{Giá thực tế và giá dự báo trên tập kiểm thử của Random Forest và MLP PyTorch.}{house-pred}
\hinhpdf{f2_bar_diab}{1}{So sánh đa tiêu chí trên tập tiểu đường.}{bar-diab}
\hinhpdf{f2_bar_house}{1}{So sánh sai số (MAE, RMSE) và hệ số xác định $R^2$ trên tập giá nhà.}{bar-house}
\hinhpdf{f2_bar_rev}{1}{So sánh Accuracy và các chỉ số macro trên tập nhận xét khách hàng.}{bar-rev}
\hinhpdf{f2_cm_diab}{1}{Ma trận nhầm lẫn trên tập kiểm thử tiểu đường: Logistic Regression, Random Forest, MLP PyTorch.}{cm-diab}
\hinhpdf{f2_cm_rev}{1}{Ma trận nhầm lẫn trên tập kiểm thử nhận xét khách hàng.}{cm-rev}

\paragraph{Thảo luận kết quả.}
\begin{itemize}
  \item \textbf{MLP tốt nhất trên dữ liệu y tế.} Cả ba bản MLP đạt F1 lớp dương quanh
    \SL{ch2_mlp_torch/diabetes/f1} (PyTorch), cao hơn Random Forest (\SL{ch2_sklearn/diabetes/rf/f1}) và vượt xa Logistic
    Regression (\SL{ch2_sklearn/diabetes/logreg/f1}). Ma trận nhầm lẫn (Hình~\ref{fig:cm-diab}) cho thấy MLP gần như không
    báo động giả (Precision \SL{ch2_mlp_torch/diabetes/precision}) nhưng vẫn bỏ sót khoảng một phần ba ca bệnh
    (Recall \SL{ch2_mlp_torch/diabetes/recall}). ROC-AUC của bốn mô hình phi tuyến chỉ chênh nhau ở chữ số thứ ba
    (\SL{ch2_sklearn/diabetes/svm/auc} -- \SL{ch2_mlp_torch/diabetes/auc}): khác biệt lớn về F1 phần nhiều đến từ
    \emph{điểm cắt} trên đường ROC chứ không từ khả năng phân biệt. Trong sàng lọc y tế, bỏ sót đắt hơn báo động giả, nên
    khi triển khai nên hạ ngưỡng của MLP xuống dưới 0,5 để đổi Precision lấy Recall.
  \item \textbf{Random Forest vẫn nhỉnh nhất trên giá nhà.} $R^2$ của Random Forest (\SL{ch2_sklearn/house/rf/r2}) cao hơn
    MLP PyTorch (\SL{ch2_mlp_torch/house/r2}) một chút, hai mô hình bám nhau trên đồ thị phân tán
    (Hình~\ref{fig:house-pred}). Cả hai đều chỉ giải thích khoảng 62\,\% phương sai: phần còn lại nằm ở thông tin dữ liệu
    không có (vị trí chính xác trong quận, dự án, chất lượng xây dựng), nên không mô hình nào vượt xa được. Cây quyết
    định chia trực tiếp không gian theo từng quận/huyện, một lợi thế tự nhiên trên dữ liệu bảng có nhiều biến phân loại.
  \item \textbf{Văn bản: MLP chính xác nhất nhưng lớp trung tính là điểm mù chung.} MLP đạt Accuracy cao nhất
    (\SL{ch2_mlp_scratch/reviews/accPct}\,\% với bản NumPy), F1 macro ngang các mô hình tuyến tính. Cả ba bản MLP quá khớp
    rất nhanh: mất mát kiểm định chạm đáy ở epoch 2 -- 3 rồi tăng (Hình~\ref{fig:loss-mlp}), vì \SL{ch2_mlp_torch/reviews/params} tham số học
    trên chưa đến 17.000 văn bản. Ma trận nhầm lẫn (Hình~\ref{fig:cm-rev}) cho thấy lớp trung tính (3 sao) bị lẫn
    sang cả hai phía: một nhận xét 3 sao thường vừa khen vừa chê.
  \item \textbf{Ba nền tảng cho kết quả tương đương.} Với cùng kiến trúc và siêu tham số, chênh lệch F1/$R^2$ giữa NumPy,
    Keras và PyTorch nằm trong khoảng dao động do khởi tạo ngẫu nhiên khác nhau (Bảng~\ref{tab:ch2-fw}). Bản NumPy tự viết
    đạt độ chính xác ngang thư viện công nghiệp, xác nhận các công thức lan truyền ngược được cài đặt đúng. Khác biệt nằm
    ở chi phí kỹ thuật: PyTorch chạy GPU, Keras có tệp lưu lớn nhất do chứa cả trạng thái bộ tối ưu.
\end{itemize}

% ---------------------------------------------------------------------
\section{Đóng gói triển khai và giao diện Web thực tế}
\label{sec:web}

\subsection{Kiến trúc triển khai hệ thống}

Toàn bộ mô hình của ba chương được phục vụ bởi một ứng dụng Flask duy nhất gồm ba tầng (Hình~\ref{fig:deploy}):
\begin{enumerate}
  \item \textbf{Tầng giao diện:} một trang HTML5/CSS3 với sáu thẻ, JavaScript gọi API bằng \texttt{fetch} và vẽ biểu đồ
    bằng SVG thuần, không phụ thuộc thư viện ngoài.
  \item \textbf{Tầng ứng dụng:} máy chủ Flask nhận yêu cầu \texttt{POST} JSON tại sáu endpoint
    \texttt{/api/diabetes}, \texttt{/api/house}, \texttt{/api/sentiment}, \texttt{/api/vision}, \texttt{/api/stock},
    \texttt{/api/churn}.
  \item \textbf{Tầng mô hình:} mọi mô hình (\texttt{.pt}, \texttt{.joblib}) được nạp lên RAM một lần lúc khởi động và suy
    luận trên CPU. Bộ tiền xử lý được dựng lại bằng chính hàm \texttt{load\_*} đã dùng khi huấn luyện nên dữ liệu người
    dùng đi qua đúng phép biến đổi mà mô hình đã học.
\end{enumerate}

\hinhtikz{f2_deploy}{Kiến trúc ba tầng của Web App phục vụ sáu phân hệ thông minh.}{deploy}

Endpoint sàng lọc tiểu đường (Đoạn mã~\ref{code:app__api_diabetes}) minh họa khuôn chung: dựng một dòng
\texttt{DataFrame} từ JSON, biến đổi bằng bộ tiền xử lý, chạy mô hình, đo thời gian và trả JSON.

\maNguon{app__api_diabetes}{Endpoint sàng lọc tiểu đường: MLP PyTorch và Random Forest cho cùng một hồ sơ}

\subsection{Giao diện thực tế các hệ thống thông minh}

\anhUI{ui_diabetes}{Phân hệ sàng lọc nguy cơ tiểu đường: người dùng nhập chỉ số lâm sàng, hệ thống trả xác suất của MLP và
Random Forest, mức nguy cơ và khuyến nghị.}{ui-diabetes}
\anhUI{ui_house}{Phân hệ định giá bất động sản: giá ước tính của Random Forest, khoảng hợp lý $\pm$ RMSE kiểm thử, dự báo
của MLP và đơn giá theo m².}{ui-house}
\anhUI{ui_sentiment}{Phân hệ phân tích cảm xúc nhận xét: xác suất ba lớp của MLP, đối chiếu Logistic Regression và các từ
đóng góp nhiều nhất.}{ui-sentiment}

Ba phân hệ trả kết quả trong vài đến vài chục mili-giây trên CPU, cho thấy các mô hình học máy cơ bản và MLP chuyển được trực tiếp
vào một quy trình nghiệp vụ thực tế.
